% Section 1: purpose and scope
\section{Purpose and scope}\label{sec:purpose}

\subsection{The problem}

Ecological inference (EI) asks how groups behaved when only group shares and
outcome shares are observed for each unit. A unit is a polling place, a
precinct, a county or a riding. For unit $i$ we see the row shares
$X_i\in\Delta^{R-1}$ (for example, how its electorate voted in election~1, or
its racial composition) and the column shares $T_i\in\Delta^{C-1}$ (for
example, how it voted in election~2). Here $\Delta^{k-1}$ is the probability
simplex in $\mathbb R^k$. We want the latent $R\times C$ table of rates
$\beta_i=[\beta_{irc}]$, where $\beta_{irc}$ is the share of row-$r$ members
of unit $i$ who chose column $c$, and above all its electorate-wide average,
the transition rate $B_{rc}$ (Section~\ref{sec:setting}).

The data link rows and columns only through the accounting identity
$T_i=\beta_i^\top X_i$. In each unit there are $RC$ unknown rates and only
$C-1$ free equations. Without further structure the rates are not
point-identified.

\subsection{Why partial identification}

The classical answers either assume enough to obtain a point (Goodman's
constancy~\citep{goodman1953}; King's truncated normal~\citep{king1997}; hierarchical R$\times$C models) or give the
widest possible range (the Duncan--Davis ``method of bounds''). The first
kind of answer can be confidently wrong when its assumption fails. The second
is always valid but often too wide to be useful.

Partial identification takes the middle road. It states an assumption, derives
the set of values of the target that are consistent with the data and that
assumption (the \emph{identified set}), and reports that set. When the
assumption is weak, the set is wide but credible. When it is strong, the set is
narrow but rests on more. The researcher can see exactly what each assumption
buys. Jiang, King, Schmaltz and Tanner~\citep{jkst2019} (JKST below) did this
for $2\times2$ tables, in the tradition of partial identification~\citep{molinari2020}, under \emph{linear contextual effects} (LCE): each
group's mean rate may depend linearly on the composition of the unit. The
present work extends that programme to general R$\times$C tables, to units
nested in districts, to sign and monotonicity beliefs from Monotone EI
\citep{egghm2025}, and to valid inference for all of these.

\subsection{The plan of the work: a ladder of assumptions}

The work organises assumptions as a ladder. Each rung adds an assumption, and
under the conditions stated in each result, its identified set lies inside the
one before it. The rungs are not all on one line: some are non-nested and are
reported side by side. Table~\ref{tab:ladder} lists them with what each adds.

\begin{longtable}{@{}L{0.17\linewidth}L{0.42\linewidth}L{0.33\linewidth}@{}}
\caption{The rungs and what each adds.}\label{tab:ladder}\\
\toprule
Rung & Assumption added (plain) & What it buys \\
\midrule\endfirsthead
\toprule Rung & Assumption added (plain) & What it buys \\ \midrule\endhead
DD (Duncan--Davis) & None beyond the accounting identity (A1). & The model-free range; the baseline everyone can accept. Districts add nothing here (Prop.~D0). \\
M2 & Each district has its own linear contextual effects. & Usually little: district bounds add up (Prop.~D5). \\
M1($\delta$) & District contextual effects differ from a shared set by at most $\delta$. & A sensitivity dial from M2 ($\delta$ large) to M1 ($\delta=0$) (Prop.~D6). \\
M1 & Shared contextual effects, district baselines (``riding dummies''). & One district with lopsided units and an extreme baseline tightens the bound for all (Prop.~D3). \\
L4($r$) & Row $r$ has no contextual effect. & Pins $(C-1)(R-1)$ unknown directions; point identification when $R=2$ (Prop.~D7(b)). \\
L5 & Constancy within each district (Goodman by district). & A point (Prop.~D7(a)); best used as a diagnostic. \\
L6 & Constancy everywhere (Goodman). & A point; historical reference. \\
M0 (pooled, beside the ladder) & One set of linear contextual effects for all units, no districts. & JKST's model in R$\times$C; non-nested with M1 (Prop.~D4). \\
M-SIGN (any linear rung) & The sign of the between or within association, or that the two reinforce each other (CR). & Large gains in simulation; a false belief gives narrow, wrong intervals (PHASE5a~\S3.4). \\
N1 & The same sign beliefs on realised unit tables, no linear model. & Bounds sharp for continuous unit tables (Prop.~M2). Not in the package. \\
N2 & Each group's mean rate is monotone in composition, any shape. & An outer LP that contains LCE with signed slopes (Prop.~N2.2). Not in the package. \\
\bottomrule
\end{longtable}

What each rung adds, in one sentence each:
\begin{itemize}[nosep]
\item \textbf{LCE (M0) over DD.} The regression of $T$ on $X$ is quadratic under
LCE and identifies the symmetric part $Q$ of the coefficient matrices. Only the
skew part $K$, of dimension $(C-1)R(R-1)/2$, is left unknown, and requiring the
implied rates to be probabilities at the extreme units bounds it.
\item \textbf{Districts (M1 and relatives) over M0.} District baselines protect
against between-district aggregation bias, and shared contextual effects carry
the information of the most informative district to all districts.
\item \textbf{Signs and reinforcement (M-SIGN) over a linear rung.} One or two
extra linear rows in $K$ can cut the bound substantially; the price is a
belief that the data cannot always check.
\item \textbf{Inference.} Every bound is computed from an estimated regression.
The plug-in bound is too narrow, so confidence intervals that treat each end
as an intersection of many estimated candidates are needed (Phase~3).
\end{itemize}

\subsection{Method map: where each method is identified and where its inference is proved}

Table~\ref{tab:methodmap} lists, for each method, the results that identify its
bound, the results that cover its confidence intervals, its status, and where the
proofs are. Appendix key: \ref{app:phase1} = PHASE1, \ref{app:phase2a} = PHASE2a,
\ref{app:phase2b} = PHASE2b, \ref{app:phase2c} = PHASE2c, \ref{app:phase3} = PHASE3,
\ref{app:phase5a} = PHASE5a, \ref{app:pkgdesign} = PACKAGE\_DESIGN.

{\footnotesize
\setlength{\tabcolsep}{3pt}
\begin{longtable}{@{}L{0.13\linewidth}L{0.24\linewidth}L{0.29\linewidth}L{0.22\linewidth}L{0.07\linewidth}@{}}
\caption{Method map.}\label{tab:methodmap}\\
\toprule Method & Identification & Inference & Status & App. \\ \midrule\endfirsthead
\toprule Method & Identification & Inference & Status & App. \\ \midrule\endhead
DD & 2b Prop.~D0; 2a Prop.~8 (cell bounds) & none needed: no estimated quantity; CIs for $B$ are intersected with DD & Proved & \ref{app:phase2b}, \ref{app:phase2a} \\
JKST ($2\times2$ M0) & 2a Prop.~4 (equivalence), Props~1--2 & P3.5 (JKST's $\mathrm{CI}_x$ is (b) with one candidate); P3.4 under tie-breaking & Proved & \ref{app:phase2a}, \ref{app:phase3} \\
M0 (pooled LCE) & 2a Lemmas~1--3, Thm~1, Props~1--2, 7, 7$'$; Cor.~1 & P3.1, P3.2, P3.4/P3.4$'$; P3.5(a), (b2), (c1); P3.6(a); P3.8. CLR = Conjecture P3.5(b3) & Identification and (a), (b2), (c1) proved in every regime; CLR algorithmic & \ref{app:phase2a}, \ref{app:phase3} \\
M2 & 2b Thm~1$'$, Thm~2, Prop.~D5 & P3.5(a), (b2) in regime R1; P3.8 in any regime; P3.6(a) & Proved (R1); R3 conjecture; R2: P3.8 only & \ref{app:phase2b}, \ref{app:phase3} \\
M1($\delta$) & 2b Thm~2(c), Prop.~D6 & as M2 & as M2 & \ref{app:phase2b}, \ref{app:phase3} \\
M1 & 2b Lemmas~D1--D2, Thm~2, Props~D3--D4 & as M2; plus P3.6(b) (M1 vs M2) and P3.6(c) (pooling), wild bootstrap; P3.7. CLR in R2: simulation finding & Proved (R1); bootstrap simulation-supported & \ref{app:phase2b}, \ref{app:phase3} \\
L4($r$), L5, L6 & 2b Prop.~D7, Thm~2(c) & P3.5 (near point identification: two-sided split; P3 \S4) & Proved & \ref{app:phase2b}, \ref{app:phase3} \\
Change across periods (M1) & 2b Prop.~D8 & via the M1 interval for the change & Proved (identification) & \ref{app:phase2b} \\
N1 (signs on realised tables) & 2c Lemma~M1, Props~M2--M4, Rem.~M2.1 & none needed (no estimated quantity) & Proved for continuous tables; integer R15 open & \ref{app:phase2c} \\
N2 (monotone means) & 2c Lemma~N2.1, Prop.~N2.2, Lemma~N2.4; Conj.~N2.2$'$ & P3.5 conditional on a prespecified degree (2c \S6); R7 open & Outer set proved; sharp set conjecture & \ref{app:phase2c} \\
M-SIGN rows & 2c Prop.~M5; 5a Thm~R3, Prop.~R2 & 5a Thm~R10(a), by interval version & Proved for pooled M0; grouped rungs R16 open & \ref{app:phase2c}, \ref{app:phase5a} \\
CR & 2c Prop.~M3; 5a Thm~R3(c)--(d) & 5a Thm~R10(b) (slab, conservative), (c) (branch union); CLR union algorithmic; R10(e) (plug-in bracket invalid) & Proved for pooled M0 & \ref{app:phase2c}, \ref{app:phase5a} \\
Ordering / M-LOC & PD Props~P-S, P-R, P-N; 2c M5(b) & 5a Thm~R10(a) as added rows (pooled M0) & P-S proved; C1 conjecture; M-LOC not built & \ref{app:pkgdesign}, \ref{app:phase2c} \\
Sensitivity $\eta$, $\delta$ & P3 Prop.~S1; 2b Prop.~D6 & CI version of the bound (P3 \S7) & Outer bound proved & \ref{app:phase3}, \ref{app:phase2b} \\
\bottomrule
\end{longtable}
}

\subsection{What this report contains and how to read it}

This report collects the theory, the package mapping and two rounds of
external referee review into one document. It is pinned to commit
\texttt{ec3644f} of branch \texttt{claude/wizardly-volta-lbh4bm} of the
repository \texttt{mdweaver/eiid}.

\begin{itemize}[nosep]
\item Sections~\ref{sec:setting}--\ref{sec:did} state each result, give its
source label in the repository (for example ``Proposition P3.4$'$ (PHASE3
\S4)''), explain it in plain terms, sketch the argument and point to the full
proof in the appendices. Where a referee changed a result, a \emph{revision
note} follows it.
\item Section~\ref{sec:evidence} collects the numerical evidence. It is
evidence, not proof.
\item Section~\ref{sec:package} maps results to the R package
\texttt{eiPartialID} 0.3.1.
\item Section~\ref{sec:status} is the status table of every result.
Section~\ref{sec:open} lists what is open.
\item Appendix~\ref{app:revisions} is the consolidated table of every referee
point, the result it changed and the change made.
\item Appendices~\ref{app:phase1}--\ref{app:pkgdesign} reproduce the
repository's theory documents verbatim (converted mechanically from Markdown).
They contain every detailed proof, with the justification tag on each step.
\end{itemize}

\paragraph{Status labels.} The referees insisted that established results,
conjectures and algorithmic procedures be kept apart. Every result carries one
of these labels:
\begin{description}[nosep,font=\normalfont\sffamily]
\item[Proved] A proof is written in the repository, in numbered steps with
justification tags, and checked numerically where possible.
\item[Proved under conditions] Proved, but only under a stated scope (for
example, the pooled model M0 only, or a fixed number of districts).
\item[Cited] The result applies a published theorem whose conditions are
checked here but whose proof is not repeated.
\item[Sketch] Only an outline of a proof exists in the repository.
\item[Conjecture] Stated precisely, supported by checks or simulation, not
proved.
\item[Algorithm] A procedure whose validity is not proved (for example, the
CLR interval as implemented). Simulation evidence cannot promote it.
\item[Open] Not settled.
\item[Proposal] Written in a Phase-1 or Phase-2 proposal document that has not
been through the review rounds (PHASE1b, PHASE2d).
\end{description}
Verification checks are tagged as in the repository: \textbf{[F]} fails on an
earlier statement and passes on the corrected one; \textbf{[V]} verifies a
corrected or new statement; \textbf{[I]} illustrates or strengthens a
statement that was already true.
